% ============================================================
%  Ricerca Operativa — shared preamble
% ============================================================

\usepackage[a4paper,
            top=25mm, bottom=25mm,
            left=25mm, right=55mm,   % wide right margin for notes
            marginparwidth=45mm,
            marginparsep=4mm]{geometry}

% --- typography & encoding ---
\usepackage[T1]{fontenc}
\usepackage[utf8]{inputenc}
\usepackage{microtype}
% Palatino for main text (newpxtext must precede amsmath)
\usepackage{newpxtext}
% Cabin as humanist sans companion to Palatino
\usepackage[scale=0.92]{cabin}

% --- math ---
\usepackage{amsmath, amssymb, amsthm, mathtools}
% Palatino-compatible math fonts (must follow amsmath)
\usepackage{newpxmath}

% --- margin notes ---
\usepackage{marginnote}
\usepackage{sidenotes}          % \sidenote{}, \marginnote{}

% --- code listings ---
\usepackage{listings}
\usepackage{xcolor}

\definecolor{codebg}{RGB}{245,245,245}
\definecolor{codecomment}{RGB}{100,100,100}
\definecolor{codekw}{RGB}{0,0,200}
\definecolor{codestr}{RGB}{163,21,21}

% accent palette — bright-sky / graphite / tomato
\definecolor{brightsky}{RGB}{48,188,237}
\definecolor{graphite}{RGB}{48,48,54}
\definecolor{ghostwhite}{RGB}{255,250,255}
\definecolor{tomato}{RGB}{252,81,48}
\definecolor{blackcolor}{RGB}{5,4,1}
\colorlet{accent}{brightsky}
\colorlet{accentgold}{tomato}
\colorlet{accentdark}{graphite}
\colorlet{accentmid}{brightsky!55!white}
\colorlet{accentlight}{brightsky!18!white}

\lstset{
  backgroundcolor=\color{codebg},
  basicstyle=\ttfamily\small,
  keywordstyle=\color{codekw}\bfseries,
  commentstyle=\color{codecomment}\itshape,
  stringstyle=\color{codestr},
  numbers=left, numberstyle=\tiny\color{gray},
  numbersep=6pt,
  breaklines=true,
  frame=single, framesep=4pt,
  captionpos=b,
  tabsize=4,
  showstringspaces=false,
}

% --- pseudocode ---
\usepackage[ruled, vlined, linesnumbered]{algorithm2e}
\SetAlgoLined
\DontPrintSemicolon

% --- graphics / figures ---
\usepackage{float}      % provides [H] float placement
\usepackage{graphicx}
\usepackage{caption}    % provides \caption*{} for unnumbered captions
\usepackage{tikz}
\usetikzlibrary{arrows.meta, positioning, shapes.geometric, decorations.pathmorphing, decorations.pathreplacing,
                matrix, fit, backgrounds, automata}

% --- booklet layout ---
\usepackage{fancyhdr}
\usepackage{titlesec}

% --- hyperlinks ---
\usepackage[colorlinks=true, linkcolor=accent,
            citecolor=graphite, urlcolor=accent]{hyperref}

% --- misc ---
\raggedbottom
\usepackage{enumitem}
\usepackage{booktabs}
\usepackage{eurosym}
\usepackage{cleveref}

% --- intermezzo box ---
\usepackage[framemethod=tikz]{mdframed}
\newmdenv[
  skipabove=\bigskipamount,
  skipbelow=\bigskipamount,
  backgroundcolor=tomato!4,
  hidealllines=true,
  innertopmargin=12pt,
  innerbottommargin=12pt,
  innerleftmargin=14pt,
  innerrightmargin=12pt,
]{intermezzoframe}
\newenvironment{intermezzo}[1]{%
  \booltrue{infullwidth}%
  \begin{adjustwidth}{-5mm}{-35mm}%
  \begin{intermezzoframe}%
  \noindent{\small\sffamily\bfseries\color{tomato!80!black}%
    \raisebox{0.5pt}{\rule{6pt}{6pt}}\enspace%
    Intermezzo\enspace---\enspace#1}%
  \par\vspace{6pt}%
}{%
  \end{intermezzoframe}%
  \end{adjustwidth}%
  \boolfalse{infullwidth}%
}

% --- formal-details box ---
% Usage: \begin{formaldetails}[optional title] ... \end{formaldetails}
% For algorithm correctness proofs use the loopproof list inside.
\usepackage{changepage}
\newmdenv[
  skipabove=\bigskipamount,
  skipbelow=\bigskipamount,
  backgroundcolor=graphite!4,
  hidealllines=true,
  innertopmargin=12pt,
  innerbottommargin=12pt,
  innerleftmargin=14pt,
  innerrightmargin=12pt,
]{formaldetailsframe}
\newenvironment{formaldetails}[1][]{%
  \booltrue{infullwidth}%
  \begin{adjustwidth}{-5mm}{-35mm}%
  \begin{formaldetailsframe}%
  \noindent{\small\sffamily\bfseries\color{graphite!80!black}%
    \raisebox{0.5pt}{\rule{6pt}{6pt}}\enspace%
    Formal details%
    \if\relax\detokenize{#1}\relax\else
      \enspace---\enspace#1%
    \fi}%
  \par\vspace{6pt}%
}{%
  \end{formaldetailsframe}%
  \end{adjustwidth}%
  \boolfalse{infullwidth}%
}

% --- loop-invariant proof structure ---
% Fields: Precondition / Loop guard / Invariant / Postcondition
%         Base case / Preservation / Postcondition validity
% Usage: \begin{loopproof} \item[Invariant] ... \end{loopproof}
% Use \lprulesep between the declaration block and the proof block.
\newlist{loopproof}{description}{1}
\setlist[loopproof]{
  font=\normalfont\small\itshape,
  leftmargin=3.2cm,
  labelwidth=3.0cm,
  labelsep=0.2cm,
  style=sameline,
  topsep=3pt,
  itemsep=4pt,
  parsep=0pt,
}
\newcommand{\lprulesep}{%
  \vspace{4pt}{\color{accentlight}\hrule height 0.4pt}\vspace{6pt}%
}

% ---- theorem environments ----
\theoremstyle{plain}
\newtheorem{theorem}{Theorem}[section]
\newtheorem{lemma}[theorem]{Lemma}
\newtheorem{corollary}[theorem]{Corollary}
\newtheorem{proposition}[theorem]{Proposition}

\theoremstyle{definition}
\newtheorem{definition}[theorem]{Definition}
\newtheorem{example}[theorem]{Example}
\newtheorem{problem}[theorem]{Problem}

\theoremstyle{remark}
\newtheorem{remark}[theorem]{Remark}
\newtheorem{note}[theorem]{Note}
\newtheorem{observation}[theorem]{Observation}
\newtheorem{exercise}[theorem]{Exercise}

% ---- plain end-of-chapter exercise environment (no box) ----
% Usage: \begin{chapex}[optional title] ... \end{chapex}
\newcounter{chapex}[chapter]
\renewcommand{\thechapex}{\arabic{chapex}}
\newenvironment{chapex}[1][]{%
  \refstepcounter{chapex}%
  \par\noindent\textbf{Exercise \thechapex%
    \if\relax\detokenize{#1}\relax\else\ (#1)\fi.}\enspace\ignorespaces
}{%
  \par\addvspace{4pt}%
}
\crefname{chapex}{Exercise}{Exercises}

% ---- theorem environment styles (colored left bar) ----
\mdfdefinestyle{thmbar}{%
  skipabove=7pt, skipbelow=5pt,
  topline=false, rightline=false, bottomline=false, leftline=true,
  linewidth=2.5pt,
  innerleftmargin=9pt, innerrightmargin=4pt,
  innertopmargin=3pt,  innerbottommargin=3pt,
}
\surroundwithmdframed[style=thmbar,
  linecolor=accent, backgroundcolor=accent!6]{theorem}
\surroundwithmdframed[style=thmbar,
  linecolor=accent, backgroundcolor=accent!6]{lemma}
\surroundwithmdframed[style=thmbar,
  linecolor=accent, backgroundcolor=accent!6]{corollary}
\surroundwithmdframed[style=thmbar,
  linecolor=accent, backgroundcolor=accent!6]{proposition}
\surroundwithmdframed[style=thmbar,
  linecolor=graphite!70!black, backgroundcolor=graphite!4]{definition}
\surroundwithmdframed[style=thmbar,
  linecolor=tomato!80!black, backgroundcolor=tomato!5]{example}
\surroundwithmdframed[style=thmbar,
  linecolor=brightsky!70!black, backgroundcolor=brightsky!4]{exercise}
\surroundwithmdframed[style=thmbar,
  linecolor=brightsky!70!black, backgroundcolor=brightsky!4]{problem}
\surroundwithmdframed[style=thmbar,
  linecolor=gray!50, backgroundcolor=gray!4]{remark}
\surroundwithmdframed[style=thmbar,
  linecolor=gray!50, backgroundcolor=gray!4]{note}
\surroundwithmdframed[style=thmbar,
  linecolor=gray!50, backgroundcolor=gray!4]{observation}

% ---- custom margin-note shorthand ----
% \aside[offset]{text} — margin note normally; inline block inside full-width environments
\usepackage{etoolbox}
\newbool{infullwidth}
\boolfalse{infullwidth}
\newcommand{\aside}[2][0pt]{%
  \ifbool{infullwidth}{%
    \begin{mdframed}[
      topline=false, rightline=false, bottomline=false,
      leftline=true, linewidth=1.5pt, linecolor=gray!35,
      innertopmargin=3pt, innerbottommargin=3pt,
      innerleftmargin=8pt, innerrightmargin=0pt,
      skipabove=6pt, skipbelow=6pt,
      backgroundcolor=white,
    ]%
    {\footnotesize\itshape\color{gray!65!black}#2}%
    \end{mdframed}%
  }{%
    \marginnote{\footnotesize\itshape #2}[#1]%
  }%
}

% ---- complexity shortcuts ----
\newcommand{\Oh}[1]{\mathcal{O}\!\left(#1\right)}
\newcommand{\Th}[1]{\Theta\!\left(#1\right)}
\newcommand{\Om}[1]{\Omega\!\left(#1\right)}

% ============================================================
%  Ricerca Operativa — OR-specific additions
% ============================================================

% --- LP formulation ---
\newcommand{\minimize}{\operatorname{minimize}}
\newcommand{\maximize}{\operatorname{maximize}}
\newcommand{\subto}{\text{subject to}}

% --- common sets ---
\newcommand{\R}{\mathbb{R}}
\newcommand{\Z}{\mathbb{Z}}
\newcommand{\N}{\mathbb{N}}

% --- simplex ---
\newcommand{\basis}{\mathcal{B}}
\newcommand{\nonbasis}{\mathcal{N}}

% --- graphs ---
\newcommand{\graph}{\mathcal{G}}
\newcommand{\vertices}{\mathcal{V}}
\newcommand{\edges}{\mathcal{E}}

% circled step numbers for algorithm figures
\newcommand{\circled}[1]{\raisebox{-0.5pt}{\textcircled{\raisebox{-0.3pt}{\scriptsize #1}}}}

% --- TikZ styles for OR ---
\tikzset{
  feasible region/.style={fill=blue!10, draw=blue!50, opacity=0.6},
  pivot cell/.style={fill=yellow!40, draw=yellow!70!black, thick},
  cut line/.style={red, dashed, very thick},
  flow edge/.style={-Stealth, thick},
  flow label/.style={midway, font=\small, fill=white, inner sep=1pt},
}

% ============================================================
%  Booklet: running headers, chapter headings, title page
% ============================================================

% --- running headers ---
\pagestyle{fancy}
\fancyhf{}
\renewcommand{\headrulewidth}{0.35pt}
\fancyhead[L]{\small\sffamily\color{gray!65!black}\nouppercase{\leftmark}}
\fancyhead[R]{\small\sffamily\color{gray!65!black}\thepage}
% chapter-opening pages: no header, page number centred at foot
\fancypagestyle{plain}{%
  \fancyhf{}%
  \renewcommand{\headrulewidth}{0pt}%
  \fancyfoot[C]{\small\sffamily\color{gray!65!black}\thepage}%
}
% use "2.  Linear Programming" in the header (not all-caps default)
\renewcommand{\chaptermark}[1]{\markboth{\thechapter.\enspace#1}{}}

% --- chapter headings ---
% CHAPTER 2          (small accent label)
% ──────────────     (thin tint rule)     [sep=6pt]
% Linear Programming (huge title)
% ═══════════════════════════════ (accent rule, after-code)
\titleformat{\chapter}[display]
  {\normalfont\sffamily\bfseries}
    {\small\color{accentmid}\MakeUppercase{\chaptertitlename}\ \thechapter}
  {6pt}
  {%
    {\color{accentlight}\rule{\linewidth}{0.5pt}}\\[6pt]%
    \Huge\color{graphite}%
  }
  [\vspace{6pt}{\color{accent}\rule{\linewidth}{2.5pt}}]
\titlespacing*{\chapter}{0pt}{28pt}{28pt}

% --- section / subsection headings ---
\titleformat{\section}
  {\large\sffamily\bfseries}
  {\color{accentmid}\thesection}{0.8em}{}
\titleformat{\subsection}
  {\normalsize\sffamily\bfseries}
  {\color{accentmid}\thesubsection}{0.7em}{}

% --- title page ---
% Full-bleed TikZ cover: accent band + thin orange strip.
% Positioned via current-page anchors so it ignores asymmetric margins.
\newcommand{\maketitlepage}{%
  \begin{titlepage}%
    \thispagestyle{empty}%
    \begin{tikzpicture}[remember picture, overlay]
      % accent top band (upper 44 % of page)
      \fill[accentdark]
        (current page.north west)
        rectangle
        ([yshift=-0.44\paperheight] current page.north east);
      % thin tomato accent strip at bottom of band
      \fill[accentgold]
        ([yshift=-0.440\paperheight] current page.north west)
        rectangle
        ([yshift=-0.453\paperheight] current page.north east);
      % main title
      \node[anchor=center,
            text=white,
            font=\fontsize{28}{34}\selectfont\sffamily\bfseries]
        at ([yshift=-0.17\paperheight] current page.north)
        {Verification and Validation Techniques};
      % subtitle
      \node[anchor=center,
            text=accentlight,
            font=\Large\sffamily\mdseries]
        at ([yshift=-0.28\paperheight] current page.north)
        {Lecture Notes};
      % author
      \node[anchor=center,
            font=\large\sffamily]
        at ([yshift=0.14\paperheight] current page.south)
        {Lucas Simonetti};
      % institution · year
      \node[anchor=center,
            text=gray!60!black,
            font=\normalsize\sffamily]
        at ([yshift=0.08\paperheight] current page.south)
        {Università degli Studi di Udine\enspace·\enspace A.A.\ 2025--2026};
    \end{tikzpicture}%
  \end{titlepage}%
}

% --- multicol package for cheatsheet ---
\usepackage{multicol}

% --- summarybox environment ---
\newmdenv[
  skipabove=\bigskipamount,
  skipbelow=\bigskipamount,
  backgroundcolor=brightsky!6,
  topline=false,
  rightline=false,
  bottomline=false,
  leftline=true,
  linewidth=3pt,
  linecolor=brightsky,
  innertopmargin=10pt,
  innerbottommargin=10pt,
  innerleftmargin=12pt,
  innerrightmargin=12pt,
]{summaryboxframe}

\newenvironment{summarybox}[1][]{%
  \begin{summaryboxframe}%
  \noindent{\small\sffamily\bfseries\color{accent}%
    \raisebox{0.5pt}{\rule{6pt}{6pt}}\enspace%
    Summary \& Key Takeaways%
    \if\relax\detokenize{#1}\relax\else
      \enspace---\enspace#1%
    \fi}%
  \par\vspace{6pt}\small\sffamily%
}{%
  \end{summaryboxframe}%
}
